% Sortierung der Literaturverzeichnisse nach Hildebrandt
\printbibheading[title={Literaturverzeichnis}]	% Kapitelueberschrift "Literatur" ausgeben.

% Allgemein wissenschaftliche Veröffentlichungen
\defbibnote{scPub}{
	Folgend werden referenzierte wissenschaftliche Publikationen aufgeführt.
}
\printbibliography[prenote=scPub, title={Wissenschaftliche Publikationen}, heading=subbibliography, keyword=WissenschaftlichePublikationen]


% Normen und Richtlinien
\defbibnote{Normen}{
	Im Folgenden werden referenzierte Normen und Richtlinien aufgeführt.
}
\printbibliography[prenote=Normen, title={Normen und Richtlinien}, heading=subbibliography, keyword=Normen]


% Publikationen des Autors
\defbibnote{ownPub}{
	Im Folgenden werden referenzierte Publikationen des Autors aufgeführt.
}
\printbibliography[prenote=ownPub, title={Publikationen des Autors}, heading=subbibliography, keyword=EigenePublikationen]


% Studentische Arbeiten
\defbibnote{StudentischeArbeiten}{
	Im Folgenden werden studentische Arbeiten aufgeführt. Diese sind mit \textsuperscript{\%} am Ende des Kurzbelegs kenntlich gemacht.
}
\printbibliography[prenote=StudentischeArbeiten, title={StudentischeArbeiten}, heading=subbibliography, keyword=StudentischeArbeiten]


% Internetquellen und Software
\defbibnote{InternetquellenSoftware}{
	Im Folgenden werden referenzierte Internetquellen und Software aufgeführt. Diese sind mit \textsuperscript{@} kenntlich gemacht.
}
\printbibliography[prenote=InternetquellenSoftware, title={Internetquellen und Software}, heading=subbibliography, keyword=InternetquellenSoftware]

% Alternativ ohne Unterteilung der einzelnen Verzeichnisse (-> davor alle vorige Zeilen auskommentieren)
% \printbibheading[title={Literaturverzeichnis}]	% Kapitelueberschrift "Literatur" ausgeben.
%\printbibliography